\documentclass[english, 11 pt, class=article, crop=false]{standalone}
\usepackage[T1]{fontenc}
%\renewcommand*\familydefault{\sfdefault} % For dyslexia-friendly text
\usepackage{lmodern} % load a font with all the characters
\usepackage{geometry}
\geometry{verbose,paperwidth=16.1 cm, paperheight=24 cm, inner=2.3cm, outer=1.8 cm, bmargin=2cm, tmargin=1.8cm}
\setlength{\parindent}{0bp}
\usepackage{import}
\usepackage[subpreambles=false]{standalone}
\usepackage{amsmath}
\usepackage{amssymb}
\usepackage{esint}
\usepackage{babel}
\usepackage{tabu}
\makeatother
\makeatletter

\usepackage{titlesec}
\usepackage{ragged2e}
\RaggedRight
\raggedbottom
\frenchspacing

% Norwegian names of figures, chapters, parts and content
\addto\captionsenglish{\renewcommand{\figurename}{Figur}}
\makeatletter
\addto\captionsenglish{\renewcommand{\chaptername}{Kapittel}}
\addto\captionsenglish{\renewcommand{\partname}{Del}}


\usepackage{graphicx}
\usepackage{float}
\usepackage{subfig}
\usepackage{placeins}
\usepackage{cancel}
\usepackage{framed}
\usepackage{wrapfig}
\usepackage[subfigure]{tocloft}
\usepackage[font=footnotesize,labelfont=sl]{caption} % Figure caption
\usepackage{bm}
\usepackage[dvipsnames, table]{xcolor}
\definecolor{shadecolor}{rgb}{0.105469, 0.613281, 1}
\colorlet{shadecolor}{Emerald!15} 
\usepackage{icomma}
\makeatother
\usepackage[many]{tcolorbox}
\usepackage{multicol}
\usepackage{stackengine}

\usepackage{esvect} %For vectors with capital letters

% For tabular
\usepackage{array}
\usepackage{multirow}
\usepackage{longtable} %breakable table

% Ligningsreferanser
\usepackage{mathtools}
\mathtoolsset{showonlyrefs}

% index
\usepackage{imakeidx}
\makeindex[title=Indeks]

%Footnote:
\usepackage[bottom, hang, flushmargin]{footmisc}
\usepackage{perpage} 
\MakePerPage{footnote}
\addtolength{\footnotesep}{2mm}
\renewcommand{\thefootnote}{\arabic{footnote}}
\renewcommand\footnoterule{\rule{\linewidth}{0.4pt}}
\renewcommand{\thempfootnote}{\arabic{mpfootnote}}

%colors
\definecolor{c1}{cmyk}{0,0.5,1,0}
\definecolor{c2}{cmyk}{1,0.25,1,0}
\definecolor{n3}{cmyk}{1,0.,1,0}
\definecolor{neg}{cmyk}{1,0.,0.,0}

% Lister med bokstavar
\usepackage[inline]{enumitem}

\newcounter{rg}
\numberwithin{rg}{chapter}
\newcommand{\reg}[2][]{\begin{tcolorbox}[boxrule=0.3 mm,arc=0mm,colback=blue!3] {\refstepcounter{rg}\phantomsection \large \textbf{\therg \;#1} \vspace{5 pt}}\newline #2  \end{tcolorbox}\vspace{-5pt}}

\newcommand\alg[1]{\begin{align} #1 \end{align}}

\newcommand\eks[2][]{\begin{tcolorbox}[boxrule=0.3 mm,arc=0mm,enhanced jigsaw,breakable,colback=green!3] {\large \textbf{Eksempel #1} \vspace{5 pt}\\} #2 \end{tcolorbox}\vspace{-5pt} }

\newcommand{\st}[1]{\begin{tcolorbox}[boxrule=0.0 mm,arc=0mm,enhanced jigsaw,breakable,colback=yellow!12]{ #1} \end{tcolorbox}}

\newcommand{\spr}[1]{\begin{tcolorbox}[boxrule=0.3 mm,arc=0mm,enhanced jigsaw,breakable,colback=yellow!7] {\large \textbf{Språkboksen} \vspace{5 pt}\\} #1 \end{tcolorbox}\vspace{-5pt} }

\newcommand{\sym}[1]{\colorbox{blue!15}{#1}}

\newcommand{\info}[2]{\begin{tcolorbox}[boxrule=0.3 mm,arc=0mm,enhanced jigsaw,breakable,colback=cyan!6] {\large \textbf{#1} \vspace{5 pt}\\} #2 \end{tcolorbox}\vspace{-5pt} }

\newcommand\algv[1]{\vspace{-11 pt}\begin{align*} #1 \end{align*}}

\newcommand{\regv}{\vspace{5pt}}
\newcommand{\mer}{\textsl{Merk}: }
\newcommand{\mers}[1]{{\footnotesize \mer #1}}
\newcommand\vsk{\vspace{11pt}}
\newcommand\vs{\vspace{-11pt}}
\newcommand\vsb{\vspace{-16pt}}
\newcommand\sv{\vsk \textbf{Svar} \vspace{4 pt}\\}
\newcommand\br{\\[5 pt]}
\newcommand{\figp}[1]{../fig/#1}
\newcommand\algvv[1]{\vs\vs\begin{align*} #1 \end{align*}}
\newcommand{\y}[1]{$ {#1} $}
\newcommand{\os}{\\[5 pt]}
\newcommand{\prbxl}[2]{
\parbox[l][][l]{#1\linewidth}{#2
	}}
\newcommand{\prbxr}[2]{\parbox[r][][l]{#1\linewidth}{
		\setlength{\abovedisplayskip}{5pt}
		\setlength{\belowdisplayskip}{5pt}	
		\setlength{\abovedisplayshortskip}{0pt}
		\setlength{\belowdisplayshortskip}{0pt} 
		\begin{shaded}
			\footnotesize	#2 \end{shaded}}}

\renewcommand{\cfttoctitlefont}{\Large\bfseries}
\setlength{\cftaftertoctitleskip}{0 pt}
\setlength{\cftbeforetoctitleskip}{0 pt}

\newcommand{\bs}{\\[3pt]}
\newcommand{\vn}{\\[6pt]}
\newcommand{\fig}[1]{\begin{figure}
		\centering
		\includegraphics[]{\figp{#1}}
\end{figure}}

\newcommand{\figc}[2]{\begin{figure}
		\centering
		\includegraphics[]{\figp{#1}}
		\caption{#2}
\end{figure}}

\newcommand{\sectionbreak}{\clearpage} % New page on each section

\newcommand{\nn}[1]{
\begin{equation}
	#1
\end{equation}
}

% Equation comments
\newcommand{\cm}[1]{\llap{\color{blue} #1}}

\newcommand\fork[2]{\begin{tcolorbox}[boxrule=0.3 mm,arc=0mm,enhanced jigsaw,breakable,colback=yellow!7] {\large \textbf{#1 (forklaring)} \vspace{5 pt}\\} #2 \end{tcolorbox}\vspace{-5pt} }
 
%colors
\newcommand{\colr}[1]{{\color{red} #1}}
\newcommand{\colb}[1]{{\color{blue} #1}}
\newcommand{\colo}[1]{{\color{orange} #1}}
\newcommand{\colc}[1]{{\color{cyan} #1}}
\definecolor{projectgreen}{cmyk}{100,0,100,0}
\newcommand{\colg}[1]{{\color{projectgreen} #1}}

% Methods
\newcommand{\metode}[2]{
	\textsl{#1} \\[-8pt]
	\rule{#2}{0.75pt}
}

%Opg
\newcommand{\abc}[1]{
	\begin{enumerate}[label=\alph*),leftmargin=18pt]
		#1
	\end{enumerate}
}
\newcommand{\abcs}[2]{
	\begin{enumerate}[label=\alph*),start=#1,leftmargin=18pt]
		#2
	\end{enumerate}
}
\newcommand{\abcn}[1]{
	\begin{enumerate}[label=\arabic*),leftmargin=18pt]
		#1
	\end{enumerate}
}
\newcommand{\abch}[1]{
	\hspace{-2pt}	\begin{enumerate*}[label=\alph*), itemjoin=\hspace{1cm}]
		#1
	\end{enumerate*}
}
\newcommand{\abchs}[2]{
	\hspace{-2pt}	\begin{enumerate*}[label=\alph*), itemjoin=\hspace{1cm}, start=#1]
		#2
	\end{enumerate*}
}

% Oppgaver
\newcommand{\opgt}{\phantomsection \addcontentsline{toc}{section}{Oppgaver} \section*{Oppgaver for kapittel \thechapter}\vs \setcounter{section}{1}}
\newcounter{opg}
\numberwithin{opg}{section}
\newcommand{\op}[1]{\vspace{15pt} \refstepcounter{opg}\large \textbf{\color{blue}\theopg} \vspace{2 pt} \label{#1} \\}
\newcommand{\ekspop}[1]{\vsk\textbf{Gruble \thechapter.#1}\vspace{2 pt} \\}
\newcommand{\nes}{\stepcounter{section}
	\setcounter{opg}{0}}
\newcommand{\opr}[1]{\vspace{3pt}\textbf{\ref{#1}}}
\newcommand{\oeks}[1]{\begin{tcolorbox}[boxrule=0.3 mm,arc=0mm,colback=white]
		\textit{Eksempel: } #1	  
\end{tcolorbox}}
\newcommand\opgeks[2][]{\begin{tcolorbox}[boxrule=0.1 mm,arc=0mm,enhanced jigsaw,breakable,colback=white] {\footnotesize \textbf{Eksempel #1} \\} \footnotesize #2 \end{tcolorbox}\vspace{-5pt} }
\newcommand{\rknut}{
Rekn ut.
}

%License
\newcommand{\lic}{\textit{Matematikken sine byggesteinar by Sindre Sogge Heggen is licensed under CC BY-NC-SA 4.0. To view a copy of this license, visit\\ 
		\net{http://creativecommons.org/licenses/by-nc-sa/4.0/}{http://creativecommons.org/licenses/by-nc-sa/4.0/}}}

%referances
\newcommand{\net}[2]{{\color{blue}\href{#1}{#2}}}
\newcommand{\hrs}[2]{\hyperref[#1]{\color{blue}\textsl{#2 \ref*{#1}}}}
\newcommand{\rref}[1]{\hrs{#1}{regel}}
\newcommand{\refkap}[1]{\hrs{#1}{kapittel}}
\newcommand{\refsec}[1]{\hrs{#1}{seksjon}}

\newcommand{\mb}{\net{https://sindrsh.github.io/FirstPrinciplesOfMath/}{MB}}


%line to seperate examples
\newcommand{\linje}{\rule{\linewidth}{1pt} }

\usepackage{datetime2}
%%\usepackage{sansmathfonts} for dyslexia-friendly math
\usepackage[]{hyperref}



\begin{document}
\chapter{Prosentrekning}
\section{Prosent av eit tal}
Lat oss no prøve å finne 20\% av 400. Dette betyr at me skal ha 20 ganger meir enn 1\% av 400 :\label{20pro}
\alg{
\text{20\% av 400}&= 20\cdot\text{1\% av 400} \\
&= 20\cdot 4\\
&= 80
}

\eks[1]{
	 Finn 50\% av 800.
	 
	 \sv
	 Me startar med å rekne ut kva 1\% av 800 er:
	 \[ 1\% \text{ av 800}=\frac{800}{100}=8 \]
	 Sidan 1\% av 800 er 8, og me skal ha 50\% av 800, gangar me 50 med 8:
	 \[ 50\cdot8 = 400 \]
	 50\% av 800 er 400.  (50\% blir alltid halvparten av talet!)
}
\eks[2]{
	Finn 70\% av 5\,000.
	
	\sv
	
	Me startar med å rekne ut kva 1\% av 5\,000 er:
	\[ 1\% \text{ av 5\,000}=\frac{5000}{100}=50 \]
	Sidan 1\% av 5\,000 er 50, og me skal ha 70\% av 5\,000, gangar me 50 med 70:
	\[ 50\cdot70 = 3500 \]
	70\% av 5\,000 er 3500.
}

\end{document}
\eks[2]{
	 Finn 2\% av 7,4. 
	
	\sv \vsb 
	\[ \frac{2\cdot 7,4}{100}=0,148 \]
}


\section{Antall prosent \textit{a} utgjør av \textit{b}}
Hva nå om vi har tallet 240 og ønsker å finne ut hvor mange prosent dette utgjør av 400?\vsk

\parbox[l][][l]{0.75\linewidth}{
	Vi vet at 1\% av 400 har verdien 4. For hver firer som går på 240 har vi derfor 1\%. Til sammen går det 60 firere på 240, dette betyr at 240 utgjør 60\% av 400.}\qquad
\parbox[r][][l]{0.15\linewidth}{\begin{shaded}%
$ \frac{240}{4}=60 $\vspace{-3pt}.\end{shaded}}
\begin{figure}[hbt]
	\centering 
	\includegraphics[scale=1]{\asym{pro2}} 
\end{figure}
En direkte utregning kan vi skrive slik:
\alg{
\text{Antall prosent 240 utgjør av 400}&=\frac{240}{\text{Verdien til 1\% av 400}} \br
&= \frac{240}{4} \br
&= 60
}

\reg[Antall prosent \boldmath $ a $ utgjør av $ b $]{
	\vs
\begin{equation}\label{autgavb}
\text{Antall prosent \textit{a} utgjør av \textit{b}}=\frac{a}{\text{1\% av }b}
\end{equation}
}
\eks[1]{Hvor mange prosent utgjør 60 av 200?

\sv
1\% av 200 er $ {\frac{200}{100}=2} $. Antall prosent 60 utgjør av 200 er derfor:
\[ \frac{60}{2}=30 \]
60 er 30\% av 200.
}
\eks[2]{Hvor mange prosent utgjør 13,1 av 12,5?
	
	\sv
	1\% av 12,5 er $ {\frac{12,5}{100}=0,125} $. Antall prosent 13,1 utgjør av 12,5 er derfor:
	\[ \frac{13,1}{0,125}=104,8 \]
	13,1 er 104,8\% av 12,5.
}
\eks[3]{
En vare kostet opprinnelig 500 kr, men prisen er satt ned til 350 kr. Hvor mange prosent avslag er gitt?

\sv
Det er gitt $ {\text{500 kr}-\text{350 kr}=\text{150 kr}} $ i avslag og 1\% av 500 er $ {\frac{500}{100}=5} $. Antall prosent 150 utgjør av 500 er:
\[ \frac{150}{5}=30  \]
Avslaget er derfor 30\%.
}
\info{Tips}{Man kan også se på prosent som et forhold man alltid uttrykker som en brøk med 100 i nevner. For eksempel er forholdet mellom 40 og 800 lik \[ \frac{40}{800}=0,05 \]
Skal vi gjøre om 0,05 til en brøk med 100 i nevener, må vi utvide brøken (Se ??):
\alg{
\frac{0,05}{1}\cdot\frac{100}{100} &= \frac{5}{100} \br
&= 5 \%
} 
40 utgjør altså 5\% av 800.\vsk

I praksis betyr dette at vi finner forholdet vi ønsker, og flytter komma to plasser til høyre. I \textsl{Eksempel 1} kunne vi skrevet:
\[ \frac{60}{200}=0,3 \]
Altså 30\%.
}
\begin{comment}
\info{Tips}{I \textsl{Eksempel 1} over fant vi at 1\% av 200 er $ {\frac{200}{100}=2} $. Å dele på 2 er derfor det samme som å dele med $ \frac{200}{100} $. Istedenfor å skrive $ \frac{60}{2} $ kunne vi altså skrevet:
\begin{flalign*}
&&  \frac{60}{\frac{100}{200}}&=60\cdot\frac{100}{200} && \llap{\color{blue}\text{Se ??}}
\end{flalign*}
Dette ser litt kronglete ut, men siden ganging med hundre bare betyr at vi flytter komma to plasser til høyre, kan vi først regne ut 
\[ \frac{60}{200}=0,3 \]
og etterpå flytte komma to plasser for å få svaret i prosent:
\[ 0,3 = 30\% \] 
}
\end{comment}
\section{Prosentfaktor}
\parbox[l][][l]{0.6\linewidth}{Vi har sett at når vi snakker om prosent, snaker vi egentlig om brøker med 100 som nevner. Disse brøkene har, som alle andre brøker, en verdi. Og for å i finne verdien til en brøk, deler vi tallet over brøkstreken med tallet under.}\qquad
\parbox[r][][l]{0.3\linewidth}{\begin{shaded}%
		Når vi deler på 100, flytter vi komma to plasser til venstre.\end{shaded}}
\newline
\vsk 
Skal vi for eksempel finne verdien av 37,9\%, får vi:
\alg{
	37,9\%&=\frac{37,9}{100} \br
	&= 0,379
}
Vi sier da at 37,7\% er skrevet som \textit{prosentfaktor}\index{prosentfaktor}.\regv
\reg[Prosentfaktor \label{profakt}]{
	Verdien av et prosenttall kalles prosentfaktoren.
}
\eks[1]{	Finn prosentfaktoren til 50\%.
	
	\sv \vs \vs
	\[ 50\% = 0,5 \]
	Prosentfaktoren er 0,5.
}
\eks[2]{	
	Finn prosentfaktoren til 3,4\%.
	
	\sv \vs \vs
	\[ 3,4\% = 0,034 \]
	Prosentfaktoren er 0,034.}
\eks[3]{	
	Finn prosentfaktoren til 123\%.
	
	\sv \vs \vs
	\[ 123\% = 1,23 \]
	Prosentfaktoren er 1,23.}\vsk
Når vi bruker prosentfaktorer sparer vi oss for å dele med 100 når vi skal finne prosenter av et tall. La oss se tilbake på regnestykket fra side \pageref{20pro}, hvor vi fant 20\% av 400. Da ganget vi 20 med 400 og delte på 100. Prosentfaktoren til 20\% er $ {\frac{20}{100}=0,2} $. Ganger vi denne med 400, får vi akkurat samme resultatet som på side \pageref{20pro}:
\alg{
\frac{20}{100}\cdot400 &= 0,2\cdot400 \\
&= 80
} 
\reg[\textit{a}\% av \textit{b}]{\vsb
\begin{equation}\label{aprob2}
a\%\text{ av }b= (\text{prosentfaktoren til }a)\cdot b
\end{equation}
}
\eks[1]{
Finn 30\% av 90.

\sv
Prosentfaktoren til 30\% er 0,3, altså får vi:
\alg{
\text{30\% av 90}&= 0,3\cdot 90 \\
&= 27
}}
\eks[2]{
Finn 24,6\% av 189,5.

\sv
Prosentfaktoren til 124,6\% er 1,246\,, altså får vi:
\alg{
	\text{124,6\% av 189,5}&= 1,246 \cdot 189,5 \\
	&=236.117
}
}
\eks[3]{
Finn 100\% av 80.

\sv
Prosentverdien til 100\% er 1, altså får vi:
\alg{
	\text{100\% av 80}&= 1\cdot 80 \\
	&= 27}
\textsl{Merk}: 100\% av et tall er alltid tallet selv!
}
\section{Vekstfaktor \label{vekstfaktor}}
I mange dagligdagse situasjoner har noe økt eller minket med en viss prosent. I en butikk kan man for eksempel komme over en skjorte som originalt koster 500 kr, men som er rabbatert med 40\%. Dette betyr at vi skal trekke ifra 40\% av originalprisen når vi skal betale.\vsk

Én måte å regne ut hva vi må betale, er å starte med å finne fratrekket:
\alg{
\text{40\% av 500}&=0.4\cdot500\\
	&=200
}
$ {500-200=300 }$, altså må vi betale 300 kr for skjorten.\vsk

Et litt annet regnestykke får vi om vi tenker på denne måten: Skal vi betale full pris, må vi betale 100\% av 500. Men får vi 40\% i rabatt, skal vi bare betale ${100\%-40\%=60\%}$ av 500:
\alg{
\text{60\% av 500}&=0.6\cdot500 \\
&= 300 
}
Svaret blir selvsagt det samme, vi må betale 300 kr for skjorten.\vsk
\end{document}
\parbox[l][][l]{0.6\linewidth}{Det er ikke alltid vi er så heldige at vi får rabatt på et produkt, ofte må vi faktisk betale et tillegg. \textit{Mervardiavgiften} er et slikt tillegg. I Norge må vi betale 25\% i merverdiavgift på mange varer.\vspace{1pt}}\qquad
\parbox[r][][l]{0.3\linewidth}{\begin{shaded}
	Merverdiavgift for-\\kortes til mva.
	\end{shaded}
	}
Det betyr at vi må betale et tillegg på 25\%, altså \\ $100\%+25\%=125\%$ av originalprisen.\vsk

\parbox[l][][l]{0.485\linewidth}{
\textsl{\color{OliveGreen}\textbf{Eksempel}}: Øreklokkene på bildet til høyre koster 999,20 kr \textsl{eksludert} mva. Men \textsl{inkludert} mva. må vi betale:
\alg{
	\text{125\% av 999,20}&=1.25\cdot999,20\\
	&= 1249
}
Altså 1249 kr.\vsk

Mer om merverdiavgiften finner du på \href{http://www.skatteetaten.no/no/Bedrift-og-organisasjon/avgifter/merverdiavgift/}{\color{blue}skatteetaten.no}.
}\quad
\parbox[r][][l]{0.55\linewidth}{
		{\vspace{4pt}
	%	\includegraphics[scale=0.4]{../fig/peltor}}}
\vsk

\prbxl{0.6}{Vi har nå sett på to eksempler: I det ene sank prisen på en vare, mens i det andre økte den. Når prisen sank med 40\%, endte vi opp med å betale 60\% av originalprisen. Vi sier da at \textit{vekstfaktoren} er 0,6. Når prisen økte med 25\%, endte vi opp med å betale 125\% av originalprisen. Da er vekstfaktoren 1,25.}\qquad
\prbxr{0.3}{Mange stusser over at ordet vekstfaktor brukes selv om en størrelse \textsl{synker}, men slik er det. Kanskje et bedre ord ville være \textit{endringsfaktor}?}
\regv
\reg[Vekstfaktor \label{vekstfakt}]{\vs
\begin{itemize}
\item Når en størrelse synker med $ a $\%, ender vi opp med $ {100\% - a\%} $ av størrelsen.

\item Når en størrelse øker med $ a $\%, ender vi opp med $ {100\% + a\%} $ av størrelsen. 

\item Verdien til $ {100\% - a\%} $ eller $ {100\% + a\%} $ kalles vekstfaktoren.
\end{itemize}
}
\eks[1]{ En vare verd 1000 kr er rabattert med 20\%.\os
\textbf{a)} Hva er vekstfaktoren?\os
\textbf{b)} Finn den nye prisen.
	
\sv 
\textbf{a)} Siden det er 20\% rabbatt må vi betale $ {100\%-20\%= 80\%} $ av originalprisen. Vekstfaktoren er derfor 0,8. \os

\textbf{b)} Den nye prisen finner vi ved å gange vekstfaktoren med originalprisen:
\[ 0,8\cdot1000  = 800 \]
Den nye prisen er altså 800 kr.
}
\eks[2]{En sjokolade koster 9,80 kr, ekskludert mva. På matvarer er det 15\% mva. \os
\textbf{a)} Hva er vekstfaktoren?	\os
\textbf{b)} Hva koster sjokoladen inkludert mva?

\sv
\textbf{a)} Med 15\% i tillegg må man betale $ {100\%+15\%= 115\%} $ av prisen eksludert mva. Vekstfaktoren er derfor 1,15.\vsk

\textbf{b)} \vs
\[ 1,15\cdot 9.90=12,25 \]
Sjokoladen koster 12,25 kr inkludert mva.
}
\section{Prosentpoeng}
Vi har akkurat sett på størrelser som økte eller minket med en viss prosent. Men hvis størrelsen selv er oppgitt i prosent, må vi holde tunga rett i munnen. La oss bruke størrelsen 10\% som et eksempel.\vsk


\prbxl{0.6}{Hvis 10\% øker med 5\%, får vi:
\alg{
\text{10\% økt med 5\%} &= 10\%\cdot1,05\\
&= 10,5\%
}
Men hvis 10\% istedenfor øker med 5 \textit{prosentpoeng}, ender vi med:
}\qquad
\prbxr{0.3}{I forrige seksjon fant vi at å øke en størrelse med 5\% er det samme som å gange størrelsen med 1,05}
\alg{
	\text{10\% økt med 5 prosentpoeng} &= 10\%+5\%\\
	&= 15\%}
10,5\% og 15\% er to helt forskjellige størrelser!\regv
\reg[Prosentpoeng]{\vs \vs
\alg{
&\text{a\% økt med \textit{b} prosentpoeng} = a\% + b\% \br
&\text{a\% minket med \textit{b} prosentpoeng} = a\% - b\%
}
}
\eks{
En dag var 5\% av elevene på en skole borte. Dagen etter var 7,5\% av elevene borte. \os
\textbf{a)} Hvor mye økte fraværet i prosentpoeng?\os
\textbf{b)} Hvor mye økte fraværet i prosent?

\sv
\textbf{a)} $ {7,5\%-5\%=2,5\%} $, derfor har fraværet økt med 2,5 prosentpoeng. \vsk

\textbf{b)} Her må vi svare på hvor mye endringen, altså 2,5\%, utgjør av 5\%. Dette er det samme som å finne hvor mye 2,5 utgjør av 5. (Se tilbake til ligning \eqref{autgavb}). 1\% av 5 er 0,05, derfor får vi:
\alg{
\text{Antall prosent 2,5 utgjør av 5}&=\frac{2,5}{0,05} \\
&= 50
}
Altså har fraværet økt med 50\%.
}
\info{Hva er egentlig forskjellen mellom prosent og prosentfaktor?}{Tenk på en skjorte som koster 200 kr. Tenk så at det er gitt 20\% rabatt på dene skjorten, altså får man $ {200\, \text{kr}\cdot0,2= 40\text{ kr}}$ i avslag. Men si at rabatten blir endret til 50\% av originalprisen, da blir avslaget $ {200\, \text{kr}\cdot0,5= 100\text{ kr}}$. \vsk

Rabatten har da gått opp fra 20\% til 50\% av originalprisen, fra 40 kr til 100 kr. En økning på 60 kr. Og nå kommer poenget: Istedenfor å spørre \textit{hvor mange prosent av originalprisen har rabatten økt?}, bruker vi ordet prosentpoeng. Det samme spørsmålet blir da \textit{hvor mange prosentpoeng har rabatten økt?} Svaret blir $ {50\%-20\% = 30\%} $, altså 30 prosentpoeng. (60 utgjør 30\% av 200).\vsk

Når vi isteden spør \textit{hvor mye har rabatten økt i prosent?}, mener vi \textit{hvor mange prosent av originalrabatten har rabatten økt?}. Dette kan vi finne på to måter:\vsk

\textsl{Metode 1}:
Originalrabatten var på 40 kr og økte med 60 kr. Hvor mange prosent 60 kr ugjør av 40 kr kan vi regne ut slik:
\begin{flalign*}
&& \frac{60}{0,4}& =150 && \llap{\color{blue}\text{1\% av 40 er 0,4}}
\end{flalign*}
Rabatten økte altså med 150\%.\vsk

\textsl{Metode 2}: \\
Økningen i prosentpoeng er 30, og startrabatten var 20\%. Hvor mange prosent 30 utgjør av 20 er:
\begin{flalign*}
&& \frac{30}{20}& =1,5 &&  \br
&& & =150\% && \llap{\color{blue}\text{Se tipset på s \pageref{protips}}}
\end{flalign*}
Rabatten økte med 150\%.
}

\end{document}

